\chapter{Conclusion}\label{sec:5}

\begin{quote}
\textbf{DRAFT --- not submission-ready.} Every number below is taken
from Chapter~\ref{sec:4} and has the same sources.
\end{quote}

\section{Summary}\label{sec:5.1}

This thesis asked whether a search agent trained with reinforcement
learning can be rewarded for \emph{grounded} answers, rather than only
correct ones, without an LLM judge and on a student budget. A
Qwen2.5-7B-Instruct agent with LoRA was trained with GRPO on HotpotQA.
It searched each question's paragraphs with BM25 and answered with
sentence-level citations, which a rule could check against HotpotQA's
gold supporting facts. Sixteen runs were trained on one GPU for about
\$49 in total: five reward and optimizer configurations on three seeds
each, plus one diagnostic run. Each was evaluated on a frozen
300-question holdout.

\textbf{RQ1: Does GRPO training improve a 7B search agent over the
un-tuned model?} Yes. Every method gains about +20 EM points over the
un-tuned model (22.3\%), and every confidence interval excludes zero.
Most of the gain comes from learning the action protocol: the un-tuned
model fails 130 of 300 episodes, while trained runs fail at most 22.

\textbf{RQ2: Does VEGR prevent the shortcuts of answer-only reward
without losing accuracy?} Yes. The two rewards reach the same accuracy
(43.6\% vs.~42.6\%, a difference of +1.0 points with CI {[}−1.9,
+3.8{]}). VEGR's grounded EM is \textbf{24.4 points higher} (CI
{[}+20.9, +28.1{]}). All three outcome-only seeds ended in a degenerate
behaviour. Two stopped citing entirely, and the third repeated the same
search in most episodes. As a result, the outcome-only agent is
\emph{less} grounded than the un-tuned model (−4.8 grounded EM points).
VEGR kept citation precision at 0.73--0.75 on every seed.

\textbf{RQ3: How do reward variants and optimizer fixes affect grounding
and search behaviour?} None improved on plain VEGR.

\begin{itemize}
\tightlist
\item
  Gated penalties (R2) and Dr.~GRPO advantages (R4) made no measurable
  difference. R2's apparent single-seed gain of +6 EM points did not
  replicate.
\item
  Doubling the coverage weight (R3) raised evidence coverage from 0.44
  to 0.57. Precision fell, and grounded EM dropped 8.0 points.
\item
  Dynamic sampling (R0-dyn, one seed) made the outcome-only agent more
  accurate (+8.0 EM points). It still cited nothing and collapsed during
  training on the same schedule as the plain outcome-only run.
\end{itemize}

The citation shortcut is therefore caused by what the reward asks for,
not only by how GRPO uses it.

\section{Contributions}\label{sec:5.2}

\begin{enumerate}
\def\labelenumi{\arabic{enumi}.}
\tightlist
\item
  \textbf{VEGR}, a judge-free, sentence-level evidence reward that also
  enforces the action protocol. Its individual terms are not new; the
  contribution is the combination and its controlled evaluation.
\item
  \textbf{Evidence that answer-only GRPO produces shortcuts} in a small,
  reproducible setting, detectable by rule. The thesis also diagnoses
  the shortcut: outcome-only runs skip about half their training steps
  for lack of reward variance, but restoring those updates does not
  prevent the shortcut.
\item
  \textbf{A controlled multi-seed comparison} with paired bootstrap
  confidence intervals and seed standard deviations. It separates reward
  design from optimizer fixes, and shows how a single-seed result can
  mislead.
\item
  \textbf{An open single-GPU pipeline.} One run costs about \$2.25.
  Splits, configurations, logs, per-episode results and the figure
  script are released in the repository.
\end{enumerate}

\section{Limitations}\label{sec:5.3}

\begin{itemize}
\tightlist
\item
  \textbf{Setting.} Retrieval is BM25 over each question's own 10
  paragraphs, on one dataset, so the search problem is far easier than
  open-web search.
\item
  \textbf{Scale.} Training uses 256 questions per run for one epoch, on
  a 7B model with LoRA.
\item
  \textbf{Statistics.} Three seeds and 300 holdout questions resolve
  large effects, such as the grounded-EM gaps. They cannot resolve EM
  differences of a few points.
\item
  \textbf{Single seed for R0-dyn.} R0-dyn has one seed and reached 250
  of its 256 target updates.
\item
  \textbf{Exploratory runs.} R3, R4 and R0-dyn were designed after
  seeing the first results.
\item
  \textbf{Loophole and weights.} The environment accepts uncited
  answers. The reward weights were not tuned.
\item
  \textbf{Coverage.} Coverage remains low. Even with VEGR, about a
  quarter of correct answers include at least one non-gold citation.
\end{itemize}

\section{Future work}\label{sec:5.4}

\begin{itemize}
\tightlist
\item
  \textbf{Close the no-citation loophole and fix the citation format.}
  Require at least one citation, and use a separator that cannot occur
  in titles (Trace 5, Section~\ref{sec:4.7}). Then test whether
  outcome-only training finds a different shortcut.
\item
  \textbf{Reward complete evidence without losing precision.} Examples
  are a coverage term that counts only gold citations, or a reward based
  on the F1 of the citation set, instead of a larger weight on coverage
  alone (R3).
\item
  \textbf{Step-level credit.} Reward individual searches by the
  information they add, in the direction of STAMP (Xu et al., 2026),
  rather than scoring only the final episode.
\item
  \textbf{Harder retrieval and more data.} Extend to an open corpus, to
  other multi-hop datasets with gold evidence such as 2WikiMultiHopQA
  and MuSiQue, and to larger training sets.
\item
  \textbf{More seeds for the diagnostic runs.} Run R0-dyn on all three
  seeds, so that the mechanism finding becomes a method comparison.
\end{itemize}
